% !TeX program = pdflatex
% ALIUS Bulletin native reconstruction source.
% Compile from the ALIUS-bulletin project root. This file does not include a pre-existing article PDF.
\providecommand{\ALIUSRootPrefix}{}

\ifdefined\ALIUSIssueBuild
\else
\documentclass[a4paper,11pt]{article}
\usepackage[margin=22mm]{geometry}
\usepackage[T1]{fontenc}
\usepackage[utf8]{inputenc}
\usepackage[hidelinks]{hyperref}
\usepackage[protrusion=true,expansion=false]{microtype}
\usepackage{parskip}
\fi
\title{Near-death experiences in the public debate: A scientific perspective}
\author{Charlotte Martial commentary on Surviving Death, 2021}
\date{ALIUS Bulletin 5}

\ifdefined\ALIUSIssueBuild
\else
\begin{document}
\fi
\maketitle
\section*{Metadata}
\textbf{Piece type:} commentary\par
\textbf{DOI:} 10.34700/0bzc-fc94\par
\textbf{Keywords:} near-death experience, public debate, consciousness, EEG\par
\section*{Abstract}
On January 6th 2021, Netflix released a new docu-series called "Surviving Death", whose first episode is dedicated to near-death experiences (NDEs). We asked ALIUS member and NDE expert Charlotte Martial (neuropsychologist and post-doctoral researcher at the Coma Science Group) to share her perspective about this episode. She watched it with great interest and shares her enthusiasm that popular media addresses this fascinating and growing area of research, which has not received the scientific and medical attention it deserves. She would like to raise several concerns about some statements and comments and address further points to move the debate about NDEs forward.\par
\section*{Extracted Text}
\begingroup
\small
\noindent ALIUS Bulletin n5 (2021) aliusresearch.org/bulletin Near-death experiences in the public debate: A scientific perspective Abstract On January 6th 2021, Netflix released a new docu-series called "Surviving Death", whose first episode is dedicated to near-death experiences (NDEs). We asked ALIUS member and NDE expert Charlotte Martial (neuropsychologist and post-doctoral researcher at the Coma Science Group) to share her perspective about this episode. She watched it with great interest and shares her enthusiasm that popular media addresses this fascinating and growing area of research, which has not received the scientific and medical attention it deserves. She would like to raise several concerns about some statements and comments and address further points to move the debate about NDEs forward. keywords: near-death experiences, public debate, consciousness, EEG Commentary on: first episode of Surviving Death (2021) Netflix NDE is a particular and touchy topic, in the sense that these experiences touch upon the notion of boundary of life and death. The existence of NDE is, to date, no longer debated in the scientific community, however, their origin is still a matter of controversy (Peinkhofer et al., 2019). To move the debate forward, it is essential that scientists consider available empirical evidence clearly and exhaustively. The significance and interpretation that experiencers give to their NDE have also been the subject of debate; they can be studied by scientists, but it is a matter of belief. We, as humans, have our own beliefs regarding death - and this is so precious! Whatever the personal beliefs regarding NDEs and what they represent, I, as a researcher, will only discuss the scientific evidence that exists to date. by Charlotte Martial Charlotte Martial cmartial@uliege.be Coma Science Group, GIGA consciousness, University of Liege, Belgium Cite as: Martial, C. (2021). Near-death experiences in the public debate: a scientific perspective by Charlotte Martial, ALIUS Bulletin, 5, 1-10, https://doi.org/10.34700/0bzc-fc94 Charlotte Martial - Near-death experience in the public debate ALIUS Bulletin n5 (2021) aliusresearch.org/bulletin 2 To this day, much of the ambiguity in the definition of NDEs mainly stems from the confusion over how to define death itself and the failure to mention the distinction between "clinical death", i.e., the cessation of heartbeat and respiration, and "brain death", i.e., the permanent cessation of functioning of the entire brain. Importantly, since the establishment of the brain death's criteria in the 1950s, no patient with these criteria has ever recovered from death. Surprisingly, at no point in this episode they mention the distinction between "clinical death" and "brain death". This might lead a viewer to envision the possibility that the experiencer may be dead, while it is not the case. They also claim that during a NDE, the brain functions are stopped. According to the best of my knowledge, there is no empirical scientific evidence of this statement. This is related to another statement that is also expressed in the episode, namely that getting EEG flatlining is necessary evidence of the complete absence of brain activity. So far, we know that current scalp-EEG technologies detect only activity common to neurons mainly in the cerebral cortex, but not deeper in the brain. Consequently, an EEG flatline might not be a reliable sign of complete brain inactivity; this limits the conclusion that can be drawn only based on EEG results. The show also discusses the case of Pamela Reynolds, which is fascinating-as all other similar testimonies. In the episode, it is claimed that she experienced her NDE when there was no brain activity. Nevertheless, the EEG data we have from this case does not permit to exclude a (neuro)physiological explanation of her NDE. In fact, we would have needed a rigorous scientific methodology to explore this case and draw any reliable conclusion about the potential (absence of) link between her NDE and underlying (neuro)physiological mechanisms. Considering the two empirical studies published by Chawla and colleagues (2009, 2017) identifying transient electrical spikes in critically ill patients just after cardiac arrest, a likely possibility would be that Pamela Reynolds has had he NDE during such electrical surges. However, this remains a hypothesis and we will never know what exactly happened to her. To move the debate forward, it is essential that scientists consider available empirical evidence clearly and exhaustively. " " Charlotte Martial - Near-death experience in the public debate ALIUS Bulletin n5 (2021) aliusresearch.org/bulletin 3 They claim that during a NDE, the brain functions are stopped. To the best of my knowledge, there is no empirical scientific evidence of this statement (...). So far, we know that current scalp-EEG technologies detect only activity common to neurons mainly in the cerebral cortex, but not deeper in the brain. " " Moreover, an important issue is that it is still unclear when NDEs are experienced exactly, that is, before, during and/or after (i.e., during recovery) the cardiac arrest for example. Indeed, the exact time of onset within the condition causing the NDE has not yet been determined. So far, based on the current scientific literature on consciousness, the most likely hypothesis is that NDEs arise when cerebral functions are still sufficiently operating. Several empirical studies also suggest the implications of various other causal agents, such as the release of endogenous neurotransmitters (see Martial et al., 2020). However, future studies are still needed to explore this issue further. In this episode, there is also a discussion about the fact that Pamela Reynolds subsequently reported elements from her environment during her operation while she was supposed to be completely unaware of external stimuli. It is important to note that, although many similar intriguing anecdotes have been reported in the literature, so far no empirical study is "methodologically strong" enough for reliably testing whether NDE experiencers did report some actual (real-life-based) events happening in the surrounding during their NDE. Taking the example of one of the most rigorous scientific studies we have so far, Parnia and co-authors (2014) claimed that one experiencer (out of a total of 330 cardiac arrest survivors) reported some elements from the surroundings during his/her cardiopulmonary resuscitation. However, it is noteworthy to mention that their protocol does not permit to exclude that those retrieved memories were not based on retrospective imaginative (re)constructions built up from other memories, expectation about the world and/or prior knowledge (see Martial et al., 2020). Charlotte Martial - Near-death experience in the public debate ALIUS Bulletin n5 (2021) aliusresearch.org/bulletin 4 it is still unclear when NDEs are experienced exactly, that is, before, during and/or after (i.e., during recovery) the cardiac arrest for example. Indeed, the exact time of onset within the condition causing the NDE has not yet been determined. " " The memory literature has repeatedly demonstrated the malleability of the human memory and that our previous memories, expectation about the world and/or prior knowledge may influence the formation of our memories. Currently, the empirical literature has not been able to confirm the accuracy of these reports. We clearly need to conduct studies with more refined methodologies to objectively examine the validity of these intriguing memories. Such methodologies could, for example, include target salient items/objects unexpected to be present in a resuscitation room and see if the patient will report it. This could be filmed by video recordings permitting to objectively examine the claims' validity of visual and auditory perception subsequently reported. I would like to stress that I do not exclude here the possibility that experiencers do report actual real-life-based events happening in the surrounding during their NDE, but I rather remind that convincing empirical evidence of this hypothesis is currently lacking. It is worth noting that if actual real-life-based events are confirmed, this will have important implications in the consciousness field and this will notably corroborate the existing empirical literature showing that unresponsiveness does not equal unconsciousness (see Sanders et al., 2012 for a review). Undoubtedly, the scientific literature on NDE has long suffered from a lack of a structured framework for analyzing the phenomenon, and especially for studying the seemingly paradoxical dissociation between the trigger event and the richness of the subjective experience, that probably occurs during a moment of brain dysfunction. Due to this lack of framework modelling this dissociation, many publications have been devoted to discussing the fact that NDEs are in support of the nonlocal consciousness theories (e.g., Carter, 2010; van Lommel, 2013; Parnia, 2007) suggesting that consciousness may not always coincide with the functioning of the brain. Charlotte Martial - Near-death experience in the public debate ALIUS Bulletin n5 (2021) aliusresearch.org/bulletin 5 it is still unclear when NDEs are experienced exactly, that is, before, during and/or after (i.e., during recovery) the cardiac arrest for example. Indeed, the exact time of onset within the condition causing the NDE has not yet been determined. " " Some proponents of this hypothesis claim that NDEs are evidence of a "dualistic" model toward the mind-brain relationship. Nonetheless, to date, convincing empirical evidence of this hypothesis is lacking. " " Some of their proponents claim that NDE is precisely what it seems to be to the individual experiencing it: an experience which can be considered as evidence of a "dualistic" model toward the mind-brain relationship. Nonetheless, to date, convincing empirical evidence of this hypothesis is lacking. It is an intriguing fact that several authors, including cardiologists (e.g., van Lommel, 2013), debating NDEs do not appear to recognize that they are states of disconnected consciousness likely underlined by brain functions, while this is currently the most likely hypothesis based on the current scientific state of the art. In reality, NDE is far from being the only example of such seemingly paradoxical dissociation and research has repeatedly shown that consciousness and behavioral responsiveness (i.e., behavioral interactions with the outside world, excluding reflex behaviors) may decouple, notably using brain imaging techniques in severe brain injury or pharmacologically induced states (e.g., Monti et al., 2010). For instance, complex patterns of cortical activity and interactions which are typically observed in awake conscious states can be observed throughout ketamine-induced unresponsiveness states after which reports of subjective experience have been recalled (Bonhomme et al., 2016; Sarasso et al., 2015). Nevertheless, it is true that NDEs are particularly intriguing due to their rich phenomenology subsequently reported. Charlotte Martial - Near-death experience in the public debate ALIUS Bulletin n5 (2021) aliusresearch.org/bulletin 6 NDE is far from being the only example of such seemingly paradoxical dissociation and research has repeatedly shown that consciousness and behavioral responsiveness (...) may decouple. " " We recently published an opinion article examining the NDE phenomenon in light of a novel framework (Martial et al., 2020), hoping that this will facilitate the development of a more nuanced description of NDEs in research, as well as in the media (see figure below). Illustration of Different States and Conditions Based on Wakefulness, Connectedness, and Internal Awareness. Reproduced from Martial et al. (2020) Research on NDEs has also long suffered from a dearth of empirical evidence regarding the NDE event itself because of its unpredictable aspect, which makes its scientific investigation extremely difficult. Nonetheless, experimental methods have been recently developed to go beyond limitations inherent in NDE research. For example, it is possible to induce resembling subjective experiences through the use of psychedelic substances (Timmermann et al., 2018), hypnosis (Martial et al., 2019) or syncope (i.e., transient cerebral hypoxia; Lempert et al., 1994a,b). Charlotte Martial - Near-death experience in the public debate ALIUS Bulletin n5 (2021) aliusresearch.org/bulletin 7 Contrary to what is said in the episode, we know that low cerebral oxygen levels can lead to very pleasant experiences - sometimes even intentionally induced by teenagers using the so-called "fainting lark" manoeuvre (Johnson et al., 1984). Notably, while investigating motor phenomena of syncope in a cohort of healthy young adults, Lempert and colleagues (1994a) were one of the first to report pleasant syncopal hallucinations. Out of 42 young adults volunteers who experienced complete syncope with falling, 25 subsequently reported visual and auditory hallucinations, such as out-of-body experiences, perceptions of lights which in some cases intensified to a glaring brightness, encountering relatives or more blurred entities, and hearing human voices (Lempert et al., 1994a,b). Seven volunteers described their experience of syncope as a negative one (notably due to disorientation), while all others had neutral or pleasant emotions - sometimes comparing them to drug-induced or meditation experiences. Some of them admitted being reluctant to "return to reality" (Lempert, 1996). In their paper, the authors qualified those memories as similar to NDEs as described by R. A. Moody in his book "Life after Life" (Moody, 1975), because of their close resemblance to subjective experiences reported after pathological severe prolonged periods of cerebral hypoxia, i.e., cardiac arrest. We should nonetheless bear in mind that these experiences are not considered as "classical" NDE, i.e., occurring in a life- threatening situation; however, these convincing experimental manipulations may help to understand the underlying mechanisms of classical NDEs. More generally than this example, the current neuroscience hypotheses are lacking (or not sufficiently described) in this episode. Finally, I would like to recall that it is too early to speculate on the universality of NDE features. Although historical descriptions of NDEs from diverse sources reveal sufficient common features which suggests a prototypical core experience that seems to be independent from cultures, societies and religions (Belanti et al., 2008; Blackmore, 1993; Greyson, 2006), large scale cross-cultural studies recruiting individuals from different cultural and religious backgrounds are currently missing. To date, publications are centrated in North America and Western Europe (Sleutjes et al., 2014). So far, it is still not clear to what extent the NDE experiencers' religiosity and cultural background influence the content of NDEs and the Charlotte Martial - Near-death experience in the public debate ALIUS Bulletin n5 (2021) aliusresearch.org/bulletin 8 interpretation of their features (Belanti et al., 2008; Blackmore, 1993). Some studies have shown a culture-related incidence of certain features, i.e., tunnel vision (Belanti et al., 2008; Kellehear, 1993; Pasricha \& Stevenson, 1986). However, most studies are case reports, thus limiting the generalizability and the conclusions that can be drawn at a cultural level. It is clear that the topic of NDE fascinates people around the world. I therefore invite media actors and scientists to be exhaustive when presenting the perspectives on NDEs, e.g., interview researchers from different laboratories testing different types of hypotheses and who recently published scientific publications in the field and have neuroscientific expertise. No matter what scientists will discover, it does not take anything away from the fact that NDE testimonies are intriguing and are identifiable as a psychological and physiological reality of clinical significance. NDE testimonies presented in the episode are, as often, moving and fascinating. I would like to use this opportunity to thank these NDE experiencers, as well as all other NDE experiencers who have shared their experience with researchers and/or journalists. Acknowledgments Dr. Charlotte Martial would like to thank ALIUS coordinators (Dr. Daniel A. Friedman, Dr. Matthieu Koroma and Dr. Raphael Milliere) for this initiative, as well as Dr. Helena Cassol and Dr. Olivia Gosseries for their astute reading and feedback. It is too early to speculate about the universality of NDE features. (...) Large scale cross-cultural studies recruiting individuals from different cultural and religious backgrounds are currently missing. " " Charlotte Martial - Near-death experience in the public debate ALIUS Bulletin n5 (2021) aliusresearch.org/bulletin 9 References Belanti, J., Perera, M., \& Jagadheesan, K. (2008). Phenomenology of near- death experiences: a cross-cultural perspective. Transcultural Psychiatry, 45(1), 121- 133. https://doi.org/10.1177/1363461507088001 Blackmore, S. (1993). Dying to live: Science and near-death experience. London: Grafton Bonhomme, V., et al. (2016) Resting-state network-specific breakdown of functional connectivity during ketamine alteration of consciousness in volunteers. Anesthesiology 125, 873-888. https://doi.org/10.1097/ALN.0000000000001275 Carter, C., ed (2010). Science and the Near-Death Experience: How Consciousness Survives Death. Simon and Schuster. Chawla, L. S., Akst, S., Junker, C., Jacobs, B., \& Seneff, M. G. (2009). Surges of electroencephalogram activity at the time of death: A case series. Journal of Palliative Medicine, 12(12), 1095-1100. https://doi.org/10.1089/jpm.2009.0159 Chawla, L. S., Terek, M., Junker, C., Akst, S., Yoon, B., Brasha-Mitchell, E., \& Seneff, M. G. (2017). Characterization of end-of-life electroencephalographic surges in critically ill patients. Death Studies, 41(6), 385-392. https://doi.org/10.1080/07481187.2017.1287138 Greyson, B. (2006). Near-death experiences and spirituality. Zygon, 41(2), 393-414. https://doi.org/10.1111/j.1467-9744.2005.00745.x Johnson, R.H., Lambie, D.G., \& Spalding, J.M.K. (1984). Syncope without heart disease. In: Johnson, R.H., Lambie, D.G., Spalding, J.M,K., editors. The interrelationships between dysfunction in the nervous and cardiovascular system. London: WB Saunders, p. 159-183 Kellehear, A. (1993). Culture, biology, and the near-death experience. A reappraisal. The Journal of Nervous and Mental Disease, 181(3), 148-156. https://doi.org/10.1097/00005053-199303000-00002 Lempert, T. (1996). Recognizing syncope: pitfalls and surprises. Journal of the Royal Society of Medicine, 89(7), 372-375. https://doi.org/10.1177/014107689608900705 Lempert, T., Bauer, M., \& Schmidt, D. (1994a). Syncope and near-death experience. Lancet, 344(8925), 829-830. https://doi.org/10.1016/s0140- 6736(94)92389-2 Lempert, T., Bauer, M., \& Schmidt, D. (1994b). Syncope: a videometric analysis of 56 episodes of transient cerebral hypoxia. Annals of Neurology, 36(2), 233- 237. https://doi.org/10.1002/ana.410360217 Martial, C., Cassol, H., Laureys, S., \& Gosseries, O. (2020). Near-death experience as a probe to explore (disconnected) consciousness. Trends in Cognitive Sciences, 24(3), 173-183. https://doi.org/10.1016/j.tics.2019.12.010 Charlotte Martial - Near-death experience in the public debate ALIUS Bulletin n5 (2021) aliusresearch.org/bulletin 10 Martial, C., Mensen, A., Charland-Verville, V., Vanhaudenhuyse, A., Rentmeister, D., Ali Bahri, M., Cassol, H., Englebert, J., Gosseries, O, Laureys, S., \& Faymonville, M-E. (2019). Neurophenomenology of near-death experience memory in hypnotic recall: a within-subject EEG study. Scientific Reports, 9(1), 1-11. https://doi.org/10.1038/s41598-019-50601-6 Monti, M.M. Vanhaudenhuyse, A., Coleman, M. R., Boly, M., Pickard, J. D., Tshibanda, L., Owen, A.M., Laureys, S. (2010) Willful modulation of brain activity in disorders of consciousness. New England Journal of Medicine, 362(7), 579-589. https://doi.org/10.1056/nejmoa0905370 Moody, R. A. (1975). Life after life. New York: Bantam Press Parnia, S. (2007). Do reports of consciousness during cardiac arrest hold the key to discovering the nature of consciousness? Medical Hypotheses, 69(4), 933-937. https://doi.org/10.1016/j.mehy.2007.01.076 Parnia, S., Spearpoint, K., De Vos, G., Fenwick, P., Goldberg, D., Yang, J., Zhu, J., Baker, K., Killingback, H. McLean, P., Wood, M. et al. (2014) AWARE - AWAreness during REsuscitation- a prospective study. Resuscitation, 85(12), 1799- 1805. https://doi.org/10.1016/j.resuscitation.2014.09.004 Pasricha, S., \& Stevenson, I. (1986). Near-death experiences in India. A preliminary report. The Journal of Nervous and Mental Disease, 174(3), 165-170. https://doi.org/10.1097/00005053-198603000-00007 Peinkhofer, C., Dreier, J., \& Kondziella, D. (2019). Semiology and mechanisms of near-death experiences. Current Neurology and Neuroscience Reports, 19(9), 1-12. https://doi.org/10.1007/s11910-019-0983-2 Sanders, R. D., Tononi, G., Laureys, S., Sleigh, J. W., \& Warner, D. S. (2012). Unresponsiveness= unconsciousness. The Journal of the American Society of Anesthesiologists, 116(4), 946-959. https://doi.org/10.1097/ALN.0b013e318249d0a7 Sarasso, S., Boly, M., Napolitani, M., Gosseries, O., Charland-Verville, V., Casarotto, S., Rosanova, M., Casali, A.G., Brichant, J.F., Boveroux, P. et al. (2015). Consciousness and complexity during unresponsiveness induced by propofol, xenon, and ketamine. Current Biology, 25(23), 3099-3105. Sleutjes, A., Moreira-Almeida, A., \& Greyson, B. (2014). Almost 40 years investigating near-death experiences. The Journal of Nervous and Mental Disease, 202(11), 833-836. https://doi.org/10.1097/nmd.0000000000000205 Timmermann, C., Roseman, L., Willimans, L., Erritzoe, D., Martial, C., Cassol, H., Laureys, S., Feilding, A., Leech, R., Nutt, D., \& Carhart-Harrit, R. (2018). DMT models the near-death experience. Frontiers in Psychology, 9, 1424. https://doi.org/10.3389/fpsyg.2018.01424 van Lommel, P. (2013) Non-local consciousness: a concept based on scientific research on near-death experiences during cardiac arrest. Journal of Consciousness Studies, 20(1-12), 47-48\par
\endgroup

\ifdefined\ALIUSIssueBuild
\clearpage
\expandafter\endinput
\fi
\end{document}
